% !TEX program = xelatex
% Bilingual document with Arabic as the main language
% and an English abstract section. Shows \LR{} for inline English
% and a full English block using \begin{english}...\end{english}.
\documentclass[12pt, a4paper]{article}

\usepackage{fontspec}
\usepackage{polyglossia}
\setmainlanguage[locale=mashriq]{arabic}
\setotherlanguage{english}

\newfontfamily\arabicfont[Script=Arabic]{Amiri}
\newfontfamily\englishfont{TeX Gyre Termes}

\begin{document}

\title{المستنداتُ ثنائيةُ اللغةِ في \LR{LaTeX}}
\author{د. نبراس أبو الذهب}
\date{2025}
\maketitle

\section{المقدمة}

يَحتاجُ الباحثُ العربيّ كثيراً إلى كتابةِ مستنداتٍ ثنائيةِ اللغة،
بحيثُ يكونُ النصُّ الرئيسيُّ بالعربيةِ معَ ملخصٍ أو مقتطفاتٍ بالإنجليزية.
تَتيحُ حزمةُ \LR{polyglossia} التبديلَ السلسَ بين اللغتين داخلَ المستندِ
الواحد. لِإدراجِ كلمةٍ أو عبارةٍ إنجليزيةٍ داخلَ النصِّ العربيّ،
نَستخدمُ \LR{\textbackslash LR\{\}} مثل: \LR{LaTeX} أو \LR{polyglossia}.

\section{الملخصُ باللغةِ العربية}

تَعالجُ هذه الورقةُ مسألةَ التنضيدِ ثنائيِّ اللغةِ في نظامِ
\LR{LaTeX}. نُقدّمُ منهجيةً متكاملةً تَجمعُ بين النصِّ العربيّ
والإنجليزيّ في مستندٍ واحدٍ معَ الحفاظِ على اتجاهِ الكتابةِ الصحيحِ
لكلِّ لغة.

\section*{Abstract (English)}

\begin{english}
This paper addresses the problem of bilingual typesetting in
\LaTeX. We present a complete methodology that combines Arabic
and English text in a single document while preserving the correct
writing direction for each language. The \texttt{polyglossia} package
with \texttt{bidi} provides the necessary infrastructure for
right-to-left layout, while inline English terms are handled via the
\texttt{\textbackslash LR\{\}} command.
\end{english}

\section{الخاتمة}

يُتيحُ هذا الإعدادُ كتابةَ أوراقٍ علميةٍ عربيةٍ معَ ملخصاتٍ إنجليزيةٍ
بسهولة، معَ الحفاظِ على جودةِ التنضيدِ في كلتا اللغتين.

\end{document}
